\documentclass{article}

% Language setting
\usepackage[spanish]{babel}
\usepackage[utf8]{inputenc}
\usepackage[T1]{fontenc}
\usepackage{graphicx}
\graphicspath{{imagenes/}}
\usepackage[colorlinks=true, allcolors=blue]{hyperref}
\usepackage[letterpaper,top=2cm,bottom=5cm,left=3cm,right=3cm,marginparwidth=1.75cm]{geometry}

% Useful packages
\usepackage{amsmath}
\usepackage{amssymb}
\usepackage{fancyhdr}
\usepackage{multicol}
\usepackage{float}
\usepackage{xcolor}
\usepackage{longtable}
\usepackage{array}
\pagestyle{fancy}

\fancyhf{}
\setlength{\headheight}{70pt}

% Marcador de sección pendiente: se muestra en rojo hasta que su responsable
% borre esta línea y escriba su contenido.
\newcommand{\pendiente}[2]{\textcolor{red}{\textit{[Pendiente --- #1: #2]}}}

% Definición del encabezado
\fancyhead[C]{%
    \begin{minipage}[c]{1.5cm}
        \includegraphics[width=1.5cm]{imagenes/uta.png}
    \end{minipage}%
    \hfill
    \begin{minipage}[c]{10cm}
        \centering
        \small {\textbf{UNIVERSIDAD TÉCNICA DE AMBATO}} \\[0.1cm]
        \scriptsize \textit{FACULTAD DE INGENIERÍA EN SISTEMAS ELECTRÓNICA E INDUSTRIAL} \\[0.1cm]
        \scriptsize \textbf{CARRERA SOFTWARE} \\[0.1cm]
        \scriptsize \textbf{CICLO ACADÉMICO: JULIO - DICIEMBRE 2026}
        \vspace{0.2cm}
    \end{minipage}%
    \hfill
    \begin{minipage}[c]{1.5cm}
        \centering
        \includegraphics[width=1.5cm]{imagenes/fisei.png}
    \end{minipage}%
}

\begin{document}
\begin{center}
    {\large \textbf{INFORME DE GUÍA PRÁCTICA}} \\
\end{center}

\section{Portada}
\begin{flushleft}
\begin{tabular}{ l p{8cm} }
    \textbf{Tema:} & Prueba práctica integradora: CaféU 2.0 \\
    \textbf{Unidad de Organización Curricular:} & Profesional \\
    \textbf{Nivel y Paralelo:} &  5 - A \\
    \textbf{Alumnos participante(s):} &
        Miranda Tarupi Luis Sebastian \newline
        Velasco Huayamave Kevin Andres \newline
        Cholota Bastidas Ariel Isai \newline
        Garcia Amores Manolo Jose \\
    \textbf{Asignatura:} & Interacción Humano-Computador \\
    \textbf{Docente:} & Caiza Caizabuano Jose Ruben \\
\end{tabular}
\end{flushleft}

\section{Informe de guía práctica}
\subsection{Objetivos}
\subsubsection{General}
Diseñar, implementar, evaluar y validar una aplicación web interactiva para la cafetería universitaria CaféU 2.0, integrando experiencia de usuario, arquitectura de información, prototipado, accesibilidad, ingeniería de software y métodos de evaluación de usabilidad.

\subsubsection{Específicos}
\begin{itemize}
    \item Analizar a la persona usuaria y su recorrido, relacionando los problemas detectados con los principios de Nielsen y Gestalt.
    \item Definir requisitos, sitemap, user flow y wireframes, y mantener su trazabilidad con las pantallas y las pruebas.
    \item Implementar catálogo, búsqueda tolerante a tildes y mayúsculas, carrito, cantidades, eliminación y confirmación, separando presentación, aplicación, dominio e infraestructura.
    \item Planificar y ejecutar la evaluación de usabilidad (heurística, recorrido cognitivo, accesibilidad y pruebas con usuarios).
    \item Triangular los resultados, priorizar el backlog, corregir al menos dos problemas y validar la versión v2.
\end{itemize}

\subsection{Modalidad}
Presencial

\subsection{Tiempo de duración}
Presenciales: 2 (120 minutos sugeridos)\\
No presenciales: 0

\subsection{Instrucciones}
Trabajo colaborativo en equipos de cuatro integrantes. Cada actividad debe entregar su evidencia en el repositorio de GitHub \url{https://github.com/Sax300-cpu/Prueba-Practica-IHC_2p}, con commits y revisiones verificables de cada integrante. No se aceptan resultados simulados o inventados como pruebas con usuarios; las métricas de muestras pequeñas son descriptivas y no se generalizan estadísticamente.

\subsection{Listado de equipos, materiales y recursos}
\textbf{Listado de equipos y materiales generales empleados en la guía práctica:}
\begin{itemize}
    \item Computadora
    \item Visual Studio Code
    \item Navegador web
    \item Git y GitHub
\end{itemize}

\textbf{TAC (Tecnologías para el Aprendizaje y Conocimiento) empleados en la guía práctica:}
\begin{itemize}
    \item[$\square$] Plataformas educativas
    \item[$\square$] Simuladores y laboratorios virtuales
    \item[$\square$] Aplicaciones educativas
    \item[$\square$] Recursos audiovisuales
    \item[$\square$] Gamificación
    \item[$\blacksquare$] Inteligencia Artificial
    \item Otros (Especifique): \underline{\hspace{5cm}}
\end{itemize}

\subsection{Actividades por desarrollar}
Desarrollar la prueba práctica integradora de CaféU 2.0 siguiendo las siete actividades indicadas: análisis de usuario y experiencia; requisitos y arquitectura de información; arquitectura de software e implementación; planificación de la evaluación; evaluación heurística y recorrido cognitivo; pilotaje y pruebas con usuarios; y triangulación, corrección y validación.

\subsection{Resultados obtenidos}
% Cada integrante edita SOLO su carpeta: secciones/<iniciales>/resultados.tex
\subsubsection{Actividad 1 y métricas --- Kevin Velasco}
% Sección de Kevin Velasco. Fuente: docs/analisis-ux.md, evaluacion/metricas.md y el historial del repositorio.
\paragraph{Organización del repositorio.} Se creó el repositorio \url{https://github.com/Sax300-cpu/Prueba-Practica-IHC_2p} con la estructura exigida (\path{docs/}, \path{prototipos/}, \path{src/}, \path{tests/}, \path{evaluacion/}, \path{evidencias/}), la rama de trabajo \texttt{develop} como rama por defecto y una rama por integrante con pull requests revisados por un compañero. Se abrieron cuatro issues, uno por integrante, con sus entregables, y se importó la versión v1 de la aplicación desde el repositorio de la materia.

\paragraph{Persona usuaria.} Daniela, 20 años, estudiante de 5.º semestre de la UTA. Usa el celular a una mano, en pasillos y patio con luz solar fuerte, durante un receso de 15 a 20 minutos. Su meta es pedir un café y algo de comer sin hacer fila y retirarlo rápido. Su perfil coincide con el de los participantes observados (estudiantes UTA ajenos al diseño y mayores de edad).

\paragraph{Touchpoints.} Se identificaron siete: (1) descubrir la app, (2) pantalla de inicio con el horario, (3) menú con catálogo y búsqueda, (4) agregar al pedido, (5) Mi pedido con cantidades y quitar, (6) confirmación y código de retiro, y (7) retiro en caja.

\paragraph{Momentos de verdad.} \textbf{MdV-1:} agregar un producto, porque si la persona duda repite el toque o abandona; dos de tres participantes no percibieron el aviso (H-01). \textbf{MdV-2:} confirmar el pedido y recibir el código, cierre de la tarea con riesgos en H-03, H-04 y H-05.

\paragraph{Curva emocional.} En una escala de 1 a 5, la emoción parte en 4 (curiosidad), baja a 3 en inicio y menú, toca su punto más bajo en 2 al agregar un producto (touchpoint 4, coincide con MdV-1), sube a 3 al revisar el pedido, a 4 al recibir el código y llega a 5 al retirar sin fila. Es una curva cualitativa construida con comentarios en voz alta de tres participantes de la evaluación previa de la materia.

\paragraph{Matriz de problemas, Nielsen y Gestalt.}
\begin{longtable}{|l|p{3.6cm}|p{3.2cm}|p{3cm}|c|}
\hline
\textbf{ID} & \textbf{Problema} & \textbf{Heurística (Nielsen)} & \textbf{Gestalt} & \textbf{Sev.} \\ \hline
H-01 & Aviso de agregado efímero y poco contrastado & N1 Visibilidad del estado & Figura-fondo & 3 \\ \hline
H-02 & Búsqueda sin resultados y sin salida & N9 Recuperación de errores & Cierre & 3 \\ \hline
H-03 & Salir de la confirmación borra el pedido & N3 Control y libertad & Continuidad & 2 \\ \hline
H-04 & Motivo del bloqueo lejos del botón & N5 Prevención de errores & Proximidad & 2 \\ \hline
H-05 & Código de retiro no copiable ni explicado & N10 Ayuda y documentación & Similitud y énfasis & 2 \\ \hline
H-06 & Contador sin unidad ni etiqueta & N2 Relación con el mundo real & Proximidad & 1 \\ \hline
H-07 & Sin indicador de pasos & N1 Visibilidad del estado & Continuidad & 1 \\ \hline
H-08 & El ítem del pedido no es tocable & N4 Consistencia y estándares & Similitud & 2 \\ \hline
H-09 & El total no se ve desde el menú & N6 Reconocer antes que recordar & Proximidad y región común & 1 \\ \hline
\end{longtable}

Los problemas de severidad 3 (H-01 y H-02) caen en el momento de verdad MdV-1 y en la búsqueda, donde la persona decide si continúa. Los hallazgos y severidades provienen de la evaluación previa de la materia (inspección de cuatro evaluadores y pruebas con tres participantes). El análisis completo está en \path{docs/analisis-ux.md}.

\paragraph{Métricas con usuarios.} Las métricas del equipo (\path{evaluacion/metricas.md}) reportan nueve ejecuciones de las tareas T1 a T3 con tres participantes (P00, P01 y P02): ocho se completaron sin ayuda (88,9\,\%), con medianas de 34 s en T1, 24 s en T2 y 43 s en T3, y un SUS promedio de 80,0 sobre 100. Con \textbf{tres participantes estas cifras son solo descriptivas} y no se generalizan estadísticamente. El consolidado no adjunta las hojas individuales de observación ni identifica el commit evaluado, por lo que su contraste con la evaluación heurística debe leerse con esa limitación.


\subsubsection{Requisitos, arquitectura de información y presentación --- Manolo García}
% Edita SOLO este archivo. Borra la linea \pendiente al escribir tu parte.
% Guia: docs/requisitos.md, docs/arquitectura-informacion.md, docs/trazabilidad.md, prototipos/wireframes y src/presentation (issue #3).
\pendiente{Manolo García}{requisitos (6 RF y 4 RNF), sitemap, user flow, 3 wireframes, matriz de trazabilidad, capa de presentación y capturas de pantalla.}


\subsubsection{Arquitectura de software, pruebas y correcciones --- Luis Miranda}
% Edita SOLO este archivo. Borra la linea \pendiente al escribir tu parte.
% Guia: src/, tests/, docs/arquitectura.md y evaluacion/backlog.md (issue #4).
\pendiente{Luis Miranda}{arquitectura en capas (presentación, aplicación, dominio, infraestructura), persistencia, pruebas automatizadas, backlog y correcciones v2}


\subsubsection{Evaluación, heurística y validación --- Ariel Cholota}
% Sección de Ariel Cholota. Fuente: evaluacion/ (tareas, heuristicas, recorrido-cognitivo, accesibilidad, metricas, validacion-v2) y evidencias/
\newcommand{\parantesdespues}[3]{%
  \begin{figure}[H]\centering
    \begin{minipage}[t]{0.3\textwidth}\centering
      \includegraphics[height=6.2cm,keepaspectratio]{#1-v1.jpg}\\[2pt]\small Antes (v1)
    \end{minipage}\hspace{1cm}
    \begin{minipage}[t]{0.3\textwidth}\centering
      \includegraphics[height=6.2cm,keepaspectratio]{#1-v2.jpg}\\[2pt]\small Después (v2)
    \end{minipage}
    \caption{#2: #3}
  \end{figure}}

\paragraph{Planificación de la evaluación.} Se definieron tres tareas críticas con consigna sin pistas, estado inicial, tiempo meta y máximo, y criterio de éxito observable (\path{evaluacion/tareas.md}): \textbf{T1} buscar un café americano sin preocuparse por tildes y agregarlo (45 s de meta, 90 s máximo); \textbf{T2} cambiar la cantidad a dos unidades y verificar el total (45 s y 90 s); \textbf{T3} quitar el producto, recuperarlo con deshacer y confirmar el pedido (60 s y 120 s). Cada tarea se acompaña de la pregunta SEQ de 1 a 7 y, al final, del cuestionario SUS. Se dispone además de plan de evaluación, consentimiento informado y hoja de observación (\path{evaluacion/}).

\paragraph{Evaluación heurística.} Se aplicaron las diez heurísticas de Nielsen y se consolidaron nueve hallazgos con severidad de 0 a 4 (H-01 a H-09, tabla de la sección de Kevin Velasco). Los dos de mayor severidad (3) son el aviso de agregado poco perceptible (H-01) y la búsqueda vacía sin salida (H-02).

\paragraph{Recorrido cognitivo.} Se analizaron las tareas T2 y T3 (\path{evaluacion/recorrido-cognitivo.md}). En T2 el control de cantidad es identificable y da retroalimentación inmediata del subtotal y el total. En T3 la secuencia quitar, deshacer y confirmar es comprensible, y la reposición conserva la cantidad exacta.

\paragraph{Accesibilidad.} Se comprobaron teclado, contraste, etiquetas y mensajes de error (\path{evaluacion/accesibilidad.md}). El flujo completo se opera con Tab, Enter y Espacio, el foco se traslada al título al cambiar de vista y los avisos se anuncian en una región de estado. Los contrastes calculados a partir de \path{src/presentation/estilos.css} son: texto principal sobre fondo 13,0:1, texto atenuado sobre fondo 5,0:1, botón rojo con texto blanco 5,4:1 y texto blanco sobre el aviso verde 4,35:1, este último ligeramente por debajo del umbral de 4,5:1 para texto normal y por mejorar.

\paragraph{Pruebas con usuarios y validación de la v2.} Los resultados de P00, P01 y P02 se consolidan en \path{evaluacion/metricas.md}, resumidos en la sección de Kevin Velasco y con la misma limitación: muestra de tres participantes, cifras descriptivas. La validación de la v2 (\path{evaluacion/validacion-v2.md}) se respalda con capturas antes y después de cinco hallazgos corregidos, tomadas con el mismo recorrido en la v1 y en la v2, y con las cinco pruebas automatizadas.

\parantesdespues{h01}{H-01}{tras agregar un producto la tarjeta indica «1 en tu pedido» y el menú muestra el total.}
\parantesdespues{h02}{H-02}{la búsqueda sin resultados nombra el filtro activo y ofrece «Limpiar búsqueda y filtros».}
\parantesdespues{h03}{H-03}{tras confirmar y volver al menú se conserva el botón «Último pedido».}
\parantesdespues{h05}{H-05}{la confirmación añade «Copiar código» y la instrucción de retiro.}
\parantesdespues{h08}{H-08}{tras quitar un producto aparece «Deshacer eliminación».}


\subsection{Habilidades blandas empleadas en la práctica}
\begin{itemize}
    \item[$\square$] Liderazgo
    \item[$\blacksquare$] Trabajo en equipo
    \item[$\square$] Comunicación asertiva
    \item[$\square$] La empatía
    \item[$\blacksquare$] Pensamiento crítico
    \item[$\square$] Flexibilidad
    \item[$\square$] La resolución de conflictos
    \item[$\square$] Adaptabilidad
    \item[$\blacksquare$] Responsabilidad
\end{itemize}

\subsection{Conclusiones}
\textbf{Kevin Velasco.} El análisis de experiencia ubicó el punto más bajo de la curva emocional en el momento de agregar un producto, y las correcciones v2 atienden ese punto con la cantidad visible en la tarjeta y el total del pedido en el menú. Trabajar con ramas y pull requests revisados dejó trazabilidad de los aportes de cada integrante. La satisfacción reportada (SUS promedio de 80/100) corresponde a solo tres participantes y es descriptiva, no concluyente.


\textbf{Manolo García.} Los tres participantes completaron la recuperación y confirmación sin ayuda. El filtro activo y la forma de descubrir el control de cantidad siguen afectando la eficiencia. El SUS medio de 80/100 describe la experiencia global en esta muestra. La persistencia se comprobó en una prueba funcional independiente mediante una recarga.


\textbf{Luis Miranda.} Separar dominio, aplicación e infraestructura permitió probar las reglas de negocio sin el navegador y cambiar la presentación sin tocar el cálculo de totales ni la persistencia. Las cinco pruebas automatizadas cubren los flujos críticos de búsqueda, cantidades, recuperación y persistencia, aunque no sustituyen las pruebas con usuarios.


\textbf{Ariel Cholota.} La planificación con consignas sin pistas y criterios de éxito observables hizo comparables las tareas T1 a T3. La evaluación heurística y el recorrido cognitivo anticiparon problemas de recuperación y visibilidad que las correcciones v2 atienden, y la auditoría de accesibilidad mostró un flujo operable por teclado con un contraste mejorable en el aviso verde. La eficacia, la eficiencia y la satisfacción se describen para tres participantes, sin generalizar.


\subsection{Recomendaciones}
\textbf{Kevin Velasco.} Repetir la evaluación con más participantes ajenos al desarrollo, conservando por sesión la hoja de observación y el consentimiento anonimizados y el commit probado, para que los resultados sean verificables. Actualizar la curva emocional con esas observaciones reales y revisar cada pull request en cruz antes de fusionarlo.


\textbf{Manolo García.} % Edita SOLO este archivo. Borra la linea \pendiente al escribir tu parte.
% Una recomendacion corta (3-5 lineas).
\pendiente{Manolo García}{recomendación para mejorar el proceso o la aplicación}


\textbf{Luis Miranda.} Agregar pruebas de la capa de presentación (por ejemplo, que el botón «Último pedido» aparezca tras confirmar), porque una refactorización de la interfaz perdió durante el proceso la etiqueta del carrito y el total del menú sin que ninguna prueba lo detectara. Generar un código de retiro distinto por pedido en lugar del código fijo heredado de la v1.


\textbf{Ariel Cholota.} % Edita SOLO este archivo. Borra la linea \pendiente al escribir tu parte.
% Una recomendacion corta (3-5 lineas).
\pendiente{Ariel Cholota}{recomendación para mejorar el proceso o la aplicación}


\subsection{Referencias}

\vspace{-3em}
\renewcommand{\refname}{}
\nocite{*}
\bibliographystyle{ieeetr}
\bibliography{referencias}

\end{document}
